% Demonstration, limitations, future work, conclusion --- ONE chapter (prd.md 3.1 item 6; SKILL
% says do not split it the way the Master's was). Label is chap:demonstration: Ch1 \ref's it five
% times under that name, never chap:conclusion.
%
% Status (author's decision, 2026-09-24): the live demonstration is PLANNED FOR THE DEFENCE, not
% performed before submission. This chapter therefore fixes the demonstration's protocol in
% advance and reports no result from it. Never write "the demonstration showed ...", a latency
% figure from live capture, or a success verdict until the demonstration has been run; when it
% has, the sentences to update are listed in .scratch/thesis-review/issues/06-ingenieur-ch1.md
% (pass-2 verifier, live-demonstration check).
%
% Ownership: every measured figure here is Ch5's (chap:validation) and must be quoted exactly as
% Ch5 gives it, with no stronger qualifier. No number that Ch1-Ch5 did not report. RQ1 and its
% answer belong to the Memoire de Master; name it in words, one sentence.
%
% What Ch1 promises this chapter (issue 06; agents/06-ingenieur-ch1/argument.md P1-P31):
%   - the protocol of the live demonstration, and what the reported runs leave undemonstrated
%     (Ch1 l.74, l.144, l.460);
%   - the scope limits of l.66-82 as limitations: simulated aircraft, kinematic simulator without
%     rotor dynamics, recorded audio, chain evaluated in parts, no airframe behaviour in wind;
%   - the clamp scope: a physics simulator or an airframe exposes no integrator to rewrite
%     (Ch1 safety problem, collision paragraph);
%   - "returns to what the measurement showed" for the safe-failure rate (Ch1 1.2, last paragraph);
%   - "where they fall short of the targets set here" (Ch1 closing): every missed criterion named.

\chapter{Demonstration, limitations and conclusion}
\label{chap:demonstration}
This chapter closes the document. Section~\ref{sec:demonstration-protocol} fixes the protocol of
the live demonstration planned for the defence, before it is run, and reports no result from it.
Section~\ref{sec:limitations} states what the measurements of Chapter~\ref{chap:validation} do not
establish, and Section~\ref{sec:future-work} the work that follows from each limitation.
Section~\ref{sec:conclusion} answers RQ2 and RQ3 against those measurements. Every figure in this
chapter is quoted from Chapter~\ref{chap:validation} and every verdict from
Table~\ref{tab:requirements-summary}; the chapter adds no measurement.

\section{Demonstration protocol}
\label{sec:demonstration-protocol}
% Author's decisions (24 Sep): two parts -- A on the kinematic backend (clamp active, 3D display
% of its state, the run of record, the only part judged), B on PyFlyt (rotor dynamics, no clamp,
% no criterion); no live-capture latency reported; the scaffold's script and four criteria.
The live-demonstration requirement of Table~\ref{tab:functional-requirements} asks for the full
pipeline to run end to end from a live microphone, with networking disabled on camera. The
demonstration is planned for the defence. Its protocol is fixed here, before the run, so that what
counts as success is decided before the outcome is known and cannot be adjusted to it. It is one
run and not an experiment: it can establish that the chain operates end to end, but it measures no
rate and no percentile, and Table~\ref{tab:requirements-summary} records the requirement as planned
rather than met.

\paragraph{What the run adds.} The demonstration is the only run that joins the parts
Chapter~\ref{chap:validation} evaluates separately. The speech experiments reported there replay
recorded audio and end on the board, the latency experiment at the state machine and the
acoustic-robustness experiment at the validator, and the formation-control experiment drives the
controller without speech. The demonstration is the first run in which live capture feeds the
pipeline, in which a spoken command travels through the state machine to the swarm controller and
the simulated vehicles, and in which the wireless link from the device to the workstation carries a
command (Section~\ref{sec:command-bus}). It is also the end-to-end run with networking disabled on
which the offline-speech-recognition and offline-operation criteria are verified; until it is made,
both remain not yet run.

\paragraph{Setup.} The Raspberry~Pi~5 runs both paths, with the thread allocation of
Table~\ref{tab:core-allocation}, from the microphone through which the golden set was recorded. The
workstation runs the consumer of the bus, the flight state machine, the swarm controller and the
simulation, as Figure~\ref{fig:architecture} places them, and the two machines are joined by their
wireless link. At the start of the run the network interfaces of both machines are displayed on
camera: every interface is disabled except that link, which is the meaning
Section~\ref{sec:operational-context} gives to networking disabled. The video records the whole run
with a visible clock, and the state machine's log, which records every rejection with the intent
rejected and the state it arrived in (Section~\ref{sec:state-machine}), is kept beside it. The clock
orders events; it does not time them. No latency is reported from the run, because a single run
yields no percentile, and every latency figure in this document remains the one measured in
Section~\ref{sec:latency-experiment}.

\paragraph{Two parts.} The run has two parts. Both use the same device, the same state machine and
the same controller with the same gains, and they differ only in the simulation backend behind the
environment interface (Section~\ref{sec:simulation-backends}). Part~A is the run of record. It flies
on the kinematic backend, the one from which every result of Section~\ref{sec:formation-control}
comes and the only one on which the separation clamp acts, and it follows the whole script of
Table~\ref{tab:demo-script}. Its vehicles are drawn in three dimensions at the positions the
kinematic backend computes, and the display adds no dynamics of its own. Part~B follows, with the
workstation's process restarted on the physics backend, and repeats steps~1, 2, 5 and~6 of the
script: the take-off, the circle, the line and the landing. Its vehicles have rigid-body and rotor
dynamics, and no separation clamp acts on them. Part~B puts the controller of
Section~\ref{sec:swarm-controller}, unretuned, under rigid-body dynamics, which otherwise only the
physics backend's tests do, and its outcome is not anticipated here. It carries no success
criterion. Separation in Part~B is
recorded rather than judged, and a close approach there falls under the limitation of
Section~\ref{sec:limitations}, not under the verdict. Omitting Part~B, if the defence leaves no time
for it, withdraws no claim.

\begin{table}[htbp]
\centering
\footnotesize
\caption[Script of the demonstration]{Script of Part~A of the demonstration, fixed before the run.
Each step is performed once, in this order. \emph{Path} is the path that carries the utterance;
\emph{State after} is the flight state once the step is complete (Table~\ref{tab:legality}); and
\emph{Expected behaviour} is what the simulated vehicles do. Part~B repeats steps 1, 2, 5 and 6 on
the physics backend.}
\label{tab:demo-script}
\begin{tabularx}{\textwidth}{@{}r >{\raggedright\arraybackslash}p{4.2cm} >{\raggedright\arraybackslash}p{1.7cm} >{\raggedright\arraybackslash}p{2.4cm} X@{}}
\toprule
Step & Utterance & Path & State after & Expected behaviour \\
\midrule
1 & ``take off to five metres'' & Parse & \texttt{TAKING\_OFF}, then \texttt{FLYING} & All five vehicles climb to 5~m \\
2 & ``form a circle with radius five metres'' & Parse & \texttt{FLYING} & The vehicles take the slots of a circle of radius 5~m \\
3 & ``move north ten metres'' & Parse & \texttt{FLYING} & The formation moves 10~m north \\
4 & ``form a line with spacing three metres'', followed at once by ``swarm hold'' & Parse, then reflex & \texttt{FLYING} & The vehicles stop lateral motion and hold altitude; the line is not formed \\
5 & ``form a line with spacing three metres'' & Parse & \texttt{FLYING} & The vehicles take the slots of a line spaced 3~m apart \\
6 & ``land'' & Parse & \texttt{LANDING}, then \texttt{LANDED} & The vehicles descend and stop on ground contact \\
7 & ``take off to five metres'', then, once airborne, ``swarm abort'' & Parse, then reflex & \texttt{FLYING}, then \texttt{ABORTED} & The vehicles climb, then descend in place, motors stopped on ground contact \\
8 & None: the reset, from the workstation & Neither & \texttt{LANDED} & None; the reset is accepted only with every vehicle on the ground and nearly at rest \\
\bottomrule
\end{tabularx}
\end{table}

\paragraph{Success criteria.} Four criteria are fixed for Part~A, and the live-demonstration
requirement is met only if all four hold. First, every step of Table~\ref{tab:demo-script} ends in
the state and the behaviour its row gives. Second, at step~4 the hold takes effect and the command
spoken before it does not. The hold may meet that command during its decode, which it then cancels,
or earlier, during transcription, in which case the ordering rule discards the command when it is
published (Section~\ref{sec:command-bus}); in either case no line formation begins. Third, no
utterance that the validator or the state machine rejects, and none that resolves to
\texttt{unknown}, produces motion: in flight a rejection becomes a hold, and otherwise nothing is
dispatched (Section~\ref{sec:validation-layers}). Fourth, the whole run is on video with its clock,
from the display of the network interfaces to the reset. A criterion that fails is reported as
failed, together with the step at which it failed.

\paragraph{What the run cannot show.} One run, by one speaker in one room, establishes that the
chain operates, not how often it does. No rate, percentile or interval follows from it, and every
criterion of Table~\ref{tab:requirements-summary} that is judged on a rate keeps the verdict
Chapter~\ref{chap:validation} gives it. The vehicles are simulated in both parts, and neither part
flies an airframe or models wind.

\paragraph{Prerequisites.} Three components the run requires are not built, and
Chapter~\ref{chap:implementation} records each where it would sit: live microphone capture
(Section~\ref{sec:audio-chain}), an entry point that starts the runtime as one process
(Section~\ref{sec:runtime}), and the link from the state machine to the swarm controller
(Section~\ref{sec:state-machine}). The link is more than a connection. The controller accepts only
a formation (Section~\ref{sec:swarm-controller}), while the script also commands a take-off, a
movement, a hold, a landing and an abort, so each of these needs a mapping onto the controller; and
the state machine's transitions at the take-off height and on ground contact, like the reset, need
the vehicles' state reported back to it. Part~A requires, in addition, the three-dimensional display
of the kinematic backend. None of the four alters a stage measured in
Chapter~\ref{chap:validation}: a capture source replaces the replay without a change to any stage
that follows it (Section~\ref{sec:audio-chain}), and the link and the display add to the chain
without changing the controller or any stage the experiments measured.

\section{Limitations}
\label{sec:limitations}
% One paragraph each; state the limitation, the evidence, and what it means for the claims.
% Quote every figure from Ch5 exactly; the pointers below are where Ch5 takes them from.
% - Simulated aircraft; the formation results come from the kinematic simulator, without rotor
%   dynamics (eval/exp4.py -> swarm/simulate.py NumpyEnv); PyFlyt is only smoke-tested. No
%   airframe behaviour in wind.
% - The clamp: written back to position only on the kinematic simulator (swarm/simulate.py), and
%   collisions are counted after it at a distance below the clamp distance, so zero collisions
%   there follows from the clamp's convergence; the intervention count is the informative figure
%   (results/exp4_formation.md). A physics simulator or an airframe exposes no position to
%   rewrite: the mechanism would have to act on velocity or thrust instead (link to future work).
% - Transfer to rigid-body dynamics not established (Ch4 sec:simulation-backends points here;
%   issue 12 V4). PyFlyt's tests fly one placement (seed 42), check only the final state, apply no
%   clamp and count no collisions. Another placement breaks both: seed 7's wedge ends at formation
%   accuracy 0.60 with a pair at 0.20 m, under the collision distance
%   (.scratch/thesis-review/agents/12-ingenieur-ch4/verifier.md A4; not in results/ -- if the
%   figure is quoted, generate it first, or state the limitation without it).
% - The speed ceiling is not applied where motion happens (Ch4 sec:swarm-controller points here;
%   issue 12 V10). The validator bounds a commanded speed at 2.0 m/s (move / set_param only);
%   both backends cap at 3 m/s, and a formation carries no speed, so formation changes can exceed
%   the ceiling. No command in the experiments set a speed. Link to the third mechanism: the clamp
%   bounds separation, not speed.
% - Recorded audio only; no live capture; live-capture latency not measured. One speaker, one
%   room, one recording session; speaker spread is the Memoire de Master's speaker-sensitivity
%   experiment (name it in words).
% - The chain evaluated in parts; no spoken command carried through to simulated flight. In the
%   latency experiment the state machine ran on the Pi over loopback, not on the workstation, and
%   the Wi-Fi hop is not in the latency figures (results/exp2_analysis.md, "the bus is loopback on
%   the Pi"). Say what that leaves out of the end-to-end figure.
% - Preemption covers only the utterance in progress (Ch3 sec:dual-path points here; issue 11
%   V2). A queued utterance takes its sequence number at dequeue, so after a spoken hold a command
%   already waiting behind it, including the spoken stopping phrase itself, compares as newer and
%   is dispatched while the swarm is FLYING. After an abort the state machine rejects it.
% - The cancellation reaches only a running decode (Ch3 sec:dual-path; issue 11 V3). A trigger
%   during transcription leaves transcription and the decode after it to run before the result is
%   discarded, so recovery can exceed the preemption-recovery criterion. The latency experiment
%   triggered only during decodes (78/78, results/exp2_analysis.md), so its recovery figure does
%   not cover this case.
% - The stop command is one unacknowledged UDP datagram (issue 11 A15). Over loopback nothing was
%   lost; over the Wi-Fi hop a lost reflex-path datagram is not retransmitted. Tie to the Wi-Fi
%   bullet above.
% - No run with networking disabled has been made; the offline-operation criterion rests on
%   construction (every model local, no external service called) until the demonstration.
% - Every criterion's verdict is in Ch5's tab:requirements-summary (generated by eval/tables.py);
%   refer to it rather than re-listing numbers, and quote a figure only as that table gives it.
% - Criteria missed: reflex latency (idle and loaded); end-to-end latency; preemption recovery;
%   clean-audio recognition and recognition in noise; the Memoire de Master's Safe failure
%   criterion (safe-failure rate far below its bound; ADR-0006 re-baselined it -- say how).
% - Criteria not demonstrated (the budget lies inside the 95% interval; declared rule for rates
%   from small event counts): keyword false accepts (point estimate above the budget) and keyword
%   false rejects on the author's voice (point estimates within it). Not "missed", not "met".
% - Not yet run: offline operation and the offline-speech-recognition check (no run with
%   networking disabled); the live demonstration is planned for the defence.
%   Per-stage attribution: which stages miss their allowance (STT, prefill) -- the answer to
%   RQ2's second clause.
% - The two error classes no layer catches (Ch1 1.3): out-of-envelope values clamped and executed;
%   well-formed, in-range, state-legal wrong commands. Tie to the safe-failure result.
% - If still unfixed at submission: the state machine judges legality on the raw intent after the
%   validator's fallback (code), and the grammar admits leading zeros (Memoire de Master) -- name
%   only if they bear on a claim here.
\TODO{draft}

\section{Future work}
\label{sec:future-work}
% Each item follows from a limitation above; no item without one.
% - Live capture and the full chain into the controller, run with networking disabled (the
%   demonstration makes this one run; a measured version would repeat it as an experiment).
% - The Wi-Fi hop measured and added to the end-to-end figure; the state machine on the workstation
%   as designed.
% - A physics simulator, then airframes: a separation mechanism that acts on commanded velocity
%   or thrust rather than on position. Cite only from references.bib (closed set); add a key
%   first if a method is named.
% - Closing the latency gap where Ch5 locates it (STT and prefill), within the reflex-path
%   membership rule, which does not change.
% - More speakers and rooms for the acoustic results.
% - A distributed transport in place of the minimal bus beyond simulation (Ch3 command bus
%   paragraph already names this as a deliberate deferral).
\TODO{draft}

\section{Conclusion}
\label{sec:conclusion}
% Unhedged, per SKILL: answer each research question directly against Ch5's figures.
% - RQ2 (pipeline latency): does the offline pipeline meet a p95 budget compatible with
%   interactive control, and where does the time go? Answer with the end-to-end and reflex
%   figures and the per-stage attribution; state plainly what was missed.
% - RQ3 (acoustic robustness and control quality): how recognition degrades from clean to 5 dB;
%   whether it degrades safely (the safe-failure result, stated without softening); whether the
%   controller holds formation (formation accuracy and collisions, with the clamp's scope).
% - The contribution: what the reflex path and its membership rule deliver as designed and as
%   measured; the rule's transferable part (Ch1 contribution section).
% - One sentence naming RQ1 as answered in the Memoire de Master, in words.
% - Close on what the demonstration at the defence will add, without anticipating its outcome.
\TODO{draft}
